\section{Introduction}
Autonomous agents leveraging Large Language Models (LLMs) present a transformative approach to decision-making by replicating human processes and workflows across various applications. These systems enhance the problem-solving capabilities of language agents by equipping them with tools and enabling collaboration with other agents, effectively breaking down complex problems into manageable components \citep{park2023generativeagentsinteractivesimulacra, havrilla2024teachinglargelanguagemodels, talebirad2023multiagentcollaborationharnessingpower, tang2024medagentslargelanguagemodels}. One prominent application of these autonomous frameworks is in the financial market—a highly complex system influenced by numerous factors, including company fundamentals, market sentiment, technical indicators, and macroeconomic events.

Traditional algorithmic trading systems often rely on quantitative models that struggle to fully capture the complex interplay of diverse factors. In contrast, LLMs excel at processing and understanding natural language data, making them particularly effective for tasks that require textual comprehension, such as analyzing news articles, financial reports, and social media sentiment. Additionally, deep learning-based trading systems often suffer from low explainability, as they rely on hidden features that drive decision-making but are difficult to interpret. Recent advancements in multi-agent LLM frameworks for finance have shown significant promise in addressing these challenges. These frameworks create explainable AI systems, where decisions are supported by evidence and transparent reasoning \citep{li2023tradinggpt, wang2024quantagentseekingholygrail, yu2024fincon}, demonstrating the potential in financial applications.

Despite their potential, most current applications of language agents in the financial and trading sectors face two significant limitations:

\textbf{Lack of Realistic Organizational Modeling:} Many frameworks fail to capture the complex interactions between agents that mimic the structure of real-world trading firms \citep{li2023tradinggpt, wang2024quantagentseekingholygrail, yu2024fincon}. Instead, they focus narrowly on specific task performance, often disconnected from the organizational workflows and established human operating procedures proven effective in trading. This limits their ability to fully replicate and benefit from real-world trading practices.

\textbf{Inefficient Communication Interfaces:} Most existing systems use natural language as the primary communication medium, typically relying on message histories or an unstructured pool of information for decision-making \citep{park2023generativeagentsinteractivesimulacra, qian2024chatdevcommunicativeagentssoftware}. This approach often results in a ``telephone effect'', where details are lost, and states become corrupted as conversations lengthen. Agents struggle to maintain context and track extended histories while filtering out irrelevant information from previous decision steps, diminishing their effectiveness in handling complex, dynamic tasks. Additionally, the unstructured pool-of-information approach lacks clear instructions, forcing logical communication and information exchange between agents to depend solely on retrieval, which disrupts the relational integrity of the data.

In this work, we address these key limitations of existing models by introducing a system that overcomes these challenges. First, our framework bridges the gap by simulating the multi-agent decision-making processes typical of professional trading teams. It incorporates specialized agents tailored to distinct aspects of trading, inspired by the organizational structure of real-world trading firms. These agents include fundamental analysts, sentiment/news analysts, technical analysts, and traders with diverse risk profiles. Bullish and bearish debaters evaluate market conditions to provide balanced recommendations, while a risk management team ensures that exposures remain within acceptable limits. Second, to enhance communication, our framework combines structured outputs for control, clarity, and reasoning with natural language dialogue to facilitate effective debate and collaboration among agents. This hybrid approach ensures both precision and flexibility in decision-making.

We validate our framework through experiments on historical financial data, comparing its performance against multiple baselines. Comprehensive evaluation metrics, including cumulative return, Sharpe ratio, and maximum drawdown, are employed to assess its overall effectiveness.